\subsection{The protected law of a designated line}
Fix a designated line and one of its occupancy stages \(N=0,1\).
Starting at the stage's entry time \(U\), suppress its next birth
through the density process
\[
 D_t=\ind_{\text{no new line birth in }(U,t]}
 \exp\left(\int_U^t\Lambda_\ell(s)\dd s\right),
\]
stopped and saved at the relevant guards or \(T\). Before \(U\) its
value is one. The birth at \(U\), if any, starts the stage and is
not counted again as its next birth. The continuous drift
\(D\Lambda_\ell\) cancels the zeroing-jump drift. Also
\[
 \int_U^t\Lambda_\ell\dd s
 =\int_{k(U)}^{k(t)}a_N(k)\dd k
 \le2\log\frac2{2-K}.
\]
Thus \(D\) is a bounded martingale and defines a probability law
conditionally on the entry history. Its upper bound is the fixed
constant \(A_0^2\), where \(A_0=2/(2-K)\). On a crossing occurring
before the next line birth, \(D\ge1\), so the original crossing
probability is at most its probability under the protected law.
In particular this comparison does not introduce the upper-density
factor \(A_0^2\). Under the protected law the next line birth has
rate zero; the external birth rates remain their original rates.

Normalize each active line weight by
\(c_p=w_p/a_N(k_\ell)\), and save this \emph{normalized} value at
positive retirement. It is not subsequently divided by a changing
\(a_N\). For a fixed cell/type subset of the line, let \(M\) be the
sum of current and saved \(c_p\). Define its entropy sum using
\[
 h(c)=c\left(1+\log\frac{2B}{c}\right),\quad h(0)=0,
 \qquad H=\sum h(c_p),
\]
again with saved terms frozen. Under the protected law the surviving
positive drift of each credit is
\(g_p=\lambda_p(1-b_N)c_p\). This follows by subtracting
\(\mathrm d\log a_N/\mathrm dt=\Lambda_\ell b_N\) from the original
ODE and deleting only the line-birth jump rates. It is
\(\lambda_pc_p/2\) at occupancy zero and zero at occupancy one.

\begin{lemma}[Credit, entropy and whole-event energy]\label{lem:entropy}
Stop further at \(M\le C_0n\). There are compensated martingales
\(Y_M,Y_H\) and a nonnegative cumulative cost \(V\) such that
\begin{align}
 M&=M(U)+D_M+Y_M,&D_M&=O(B),\label{eq:creditdecomp}\\
 H&=H(U)+D_H-V+Y_H,&D_H&=O(B\log n),\label{eq:entropydecomp}
\end{align}
and, before the virtual guard,
\begin{align}
 \langle Y_M\rangle&\le C n^\gamma V,
 &|\Delta Y_M|&\le C n^\gamma,\label{eq:bracketM}\\
 \langle Y_H\rangle&\le C n^\gamma(\log n)^2V,
 &|\Delta Y_H|&\le C n^\gamma\log n.\label{eq:bracketH}
\end{align}
The constants are independent of the admissible entry history.
\end{lemma}
\begin{proof}
For a multiplier \(J\in\{0,1/2,1\}\), direct algebra gives
\[
 h'(c)c(1-J)+h(Jc)-h(c)=-c(J\log J-J+1).
\]
The cost is \(c\) for deletion, \(c(1-\log2)/2\) for halving,
and zero otherwise. It is comparable by fixed positive constants
to the credit decrease \(c(1-J)\). Define \(V\) to be the rate
integral of the \emph{sum of these costs for each complete birth
event}. Compensating gives \eqref{eq:creditdecomp} and
\eqref{eq:entropydecomp}. Since \(\lambda_p\le B/n\) and the
credit sum is capped by \(C_0n\), the first drift is \(O(B)\).
The active credits have the positive lower bound supplied by
Lemma~\ref{lem:product} divided by \(A_0\), so
\(h'(c)=\log(2B/c)=O(\log n)\) on them. This gives the entropy
drift estimate.

In one external birth event, the ordinary-line deletion of normalized
mass is at most \(M_n\), and effects through other designated
directions touch at most \(g\) points, each of credit at most \(2B\).
Its total credit decrease \(\Delta\) is therefore at most
\(M_n+2gB=O(n^\gamma)\). The square of the \emph{whole} decrease
satisfies \(\Delta^2\le C n^\gamma\Delta\).
The absolute entropy decrease is at most \(C\log n\) times this
credit decrease, including deletion to zero. Comparing the decrease
to the event's total cost and integrating its birth rate proves
\eqref{eq:bracketM}--\eqref{eq:bracketH}. These estimates retain
all cross terms of simultaneous changes. Summing the squares of
individual point changes would not prove the required bracket bound.
The birth that causes a global stop is included: its decrease is
controlled by the virtual-guard bound in the state immediately
before that birth. Positive saving itself changes neither the
credit ledger nor its entropy.
\end{proof}

\begin{lemma}[Exponential crossing estimate]\label{lem:exp}
Let \(Y\) be a compensated birth martingale with \(Y_0=0\), jumps
bounded in absolute value by \(b\), and continuous predictable
quadratic variation \(Q\). For \(x,v>0\),
\begin{equation}\label{eq:freedman}
 \Prob(\exists t\le T:Y_t\ge x,\ Q_t\le v)
 \le\exp\left(-\frac{x^2}{2(v+bx/3)}\right).
\end{equation}
\end{lemma}
\begin{proof}
For \(0<\theta b<3\) and \(|u|\le b\), the Taylor series and
\(k!\ge2\cdot3^{k-2}\) for \(k\ge2\) give
\[
 e^{\theta u}-1-\theta u
 \le\psi u^2,\qquad \psi=\frac{\theta^2}{2(1-\theta b/3)}.
\]
Stop at the first \(Y\ge x\), capped at \(T\). The stopped
exponential \(\exp(\theta Y-\psi Q)\) is bounded by
\(\exp(\theta(x+b))\). Its compensated drift is nonpositive by the
displayed inequality, so its expectation is at most one.
Since \(Q\) is continuous and nondecreasing, a crossing with
\(Q\le v\) exists precisely when the first crossing has \(Q\le v\).
On this event the terminal stopped exponential is at least
\(\exp(\theta x-\psi v)\). Markov gives
\(\exp(-\theta x+\psi v)\); take \(\theta=x/(v+bx/3)\).
This proof uses a bounded stopped exponential, rather than assuming
that an unlocalized exponential local martingale has mean one.
\end{proof}

\begin{proposition}[Clock stops are rare]\label{prop:clock}
For each fixed positive tolerance and admissible fixed data, the
failure probability of the simultaneous designated-line credit
bound is at most
\[
 Cn\exp\left(-c\frac{n^{1-\gamma}}{\log n}\right).
\]
In particular, the probability of a designated clock reaching \(K\)
is smaller than every fixed inverse power of \(n\).
\end{proposition}
\begin{proof}
Stop the cumulative cost at \(C_1n\log n\). Since \(H\ge0\),
\(H(U)=O(n\log n)\), and \(D_H=o(n\log n)\), reaching this cost
cap, for sufficiently large \(C_1\), forces an upward fluctuation
of \(Y_H\) of order \(n\log n\). Apply Lemma~\ref{lem:exp} with
\(v=O(n^{1+\gamma}(\log n)^3)\) and
\(b=O(n^\gamma\log n)\). An increase of \(M\) by a fixed multiple
of \(n\) similarly forces a fluctuation of \(Y_M\) of order \(n\),
with \(v=O(n^{1+\gamma}\log n)\) and \(b=O(n^\gamma)\).
Both exponents have order \(n^{1-\gamma}/\log n\).

At the first birth on the line,
\((w/2)/a_1=w/a_0\); the selected point is removed. The saved ledger
is unchanged and the moving normalized credit cannot increase.
At the second birth, all remaining moving credit disappears. To
treat the occupancy-one stage, condition on its actual entry history
and the preceding good credit bound. At time zero the ledger of a
cell/type piece equals its lattice-point count \(L_E\le n\).
On a good occupancy-zero history it is at most
\(L_E+\varepsilon n\), and the first line birth cannot increase it.
The lower and upper credit bounds from Lemma~\ref{lem:product}
give \(h(c)\le C(\log n)c\), including saved positive credits.
Consequently both initial credit and entropy bounds hold uniformly
at the occupancy-one entry time. Apply the same conditional bound
to every such history and average it. The union is only over the
two stages, not over possible histories or starting times.

There are \(O(n)\) designated lines and a fixed number of cell/type
pieces. The estimates control whole time intervals, so no separate
union over phases is needed. Summing their crossing probabilities preserves
the displayed bound. On the resulting good event, integrate the
moving normalized credit with the masks to obtain
\[
 k_\ell\le\text{deterministic line capacity}
             +C_{\mathrm{data}}T\varepsilon+O(n^{-1}).
\]
Choose \(\varepsilon>0\) first so that the second term is smaller
than half of \(\kappa_H\), then choose \(n\) large. The capacity
condition \eqref{eq:heavyinput} precludes hitting \(K\).
Positive retirement need not be rare for this argument: its credit
was retained in the ledger throughout.
\end{proof}

\subsection{Explicit clock constants for the verified input}

For the input used in the application,
\[
 g=16,\quad K=255/128,\quad A_0=256,\quad G=256,
 \quad T=\frac{30088384743729}{25000000000}<2^{11}.
\]
Write \(C_{\rm line}=4G=1024\). The verified designated-line
capacity bound has margin
\[
 \kappa_H=\frac{69405753}{80000000000},\qquad
 \varepsilon_H=\frac{\kappa_H}{8TC_{\rm line}}
       =\frac{115676255}{1314581588376682496}>2^{-34}.
\]
The following explicit estimate is not needed to compute a final
threshold for all five stopping causes. It records how the
history-uniform clock argument applies to this particular input.

\begin{corollary}[A quantitative clock budget]\label{cor:clockexplicit}
For all sufficiently large \(n\),
\begin{equation}\label{eq:clockexplicit}
 \sup_p\Prob^{(p)}(\text{clock is an early cause})
 \le221184\,n^{1+d}
 \exp\left(-\frac{\varepsilon_H^2 n^{1-\gamma}}
                     {2^{18}\log n}\right).
\end{equation}
In fact, with \(d=19999/20000\) and \(\gamma=39999/40000\),
\begin{equation}\label{eq:clockthreshold}
 \log n\ge10^7\quad\Longrightarrow\quad
 \sup_p\Prob^{(p)}(\text{clock is an early cause})<2\cdot10^{-6}.
\end{equation}
\end{corollary}
\begin{proof}
Assume first \(n^{1-d}\ge30T\). The self-correction in
\eqref{eq:product} has absolute value at most \(15TB/n\le1/2\),
so \(w_p\ge2^{-17}\) while active. Every positive normalized
credit, including a saved credit, therefore satisfies
\[
 2^{-25}\le c_p\le2B.
\]
For \(L_n=32\log n\) and \(n\ge2\), this implies
\(h(c)\le L_nc\) and \(0\le h'(c)\le L_n\) on positive active
credits. Stop temporarily while \(M\le3n\). The drift terms in
Lemma~\ref{lem:entropy} obey
\[
 0\le D_M\le\tfrac32TB,\qquad 0\le D_H\le L_nD_M.
\]
For a complete external birth let \(v\) be its entropy cost and
\(\Delta\) its credit loss. Since \(\log2<3/4\), deletion and
halving both satisfy \(\Delta/4\le v\le\Delta\). The virtual
guard and at most 16 other designated-line intersections give
\(\Delta\le n^\gamma+32B\le33n^\gamma\). Hence
\begin{align*}
 \langle Y_M\rangle&\le132n^\gamma V,
 &|\Delta Y_M|&\le33n^\gamma,\\
 \langle Y_H\rangle&\le132n^\gamma L_n^2V,
 &|\Delta Y_H|&\le33n^\gamma L_n.
\end{align*}
The bound on an entropy jump follows directly from
\(h(c)-h(Jc)\le L_nc(1-J)\), including \(J=0\);
it does not assume a bounded derivative at zero.

Consider a stage starting with \(M(U)\le2n\), and any
\(0<\varepsilon\le1\) for which
\(n^{1-d}\ge3T/\varepsilon\). Stop at the first of a credit
increase \(\varepsilon n\), cost \(V=4L_nn\), an original guard,
and \(T\). Credit increases and \(V\) are continuous, so before
this stop \(M\le3n\), and \(D_M\le\varepsilon n/2\).
If the cost boundary is reached first, \(H\ge0\),
\(H(U)\le2L_nn\) and \(D_H\le L_n\varepsilon n/2\) force
\(Y_H\ge L_nn\). Lemma~\ref{lem:exp}, with
\(v=528n^{1+\gamma}L_n^3\) and \(b=33n^\gamma L_n\), bounds
its probability by
\(\exp[-n^{1-\gamma}/(1078L_n)]\).
If the credit boundary is reached first, \(Y_M\ge\varepsilon n/2\)
and \(\langle Y_M\rangle\le528L_nn^{1+\gamma}\), giving
\(\exp[-\varepsilon^2n^{1-\gamma}/(4268L_n)]\).
As \(4268\cdot32<2^{18}\), their sum is at most
\[
 2\exp\left(-\frac{\varepsilon^2n^{1-\gamma}}
                         {2^{18}\log n}\right).
\]
This bound is conditional on the entry history and its constants
are independent of that history.

At the first line birth the ledger cannot increase. Applying the
last estimate first at occupancy zero and then conditionally at
occupancy one shows that
\(M_E(t)\le L_E+2\varepsilon n\) fails with at most twice the
displayed probability. The 16 primitive unoriented directions of
height at most three satisfy
\(\sum_{v\in\Gamma}(|v_1|+|v_2|)=54\), so they determine at
most \(54n\) lines. Each line has at most 1024 nonempty cell/type
pieces, by Lemma~\ref{lem:cells} and the corresponding axis count.
Thus the simultaneous ledger bound fails under the original law
with probability at most
\[
 221184\,n\exp\left(-\frac{\varepsilon^2n^{1-\gamma}}
                                  {2^{18}\log n}\right).
\]

On its complement, integration against the masks gives
\[
 k_\ell(t)\le\frac1n\sum_{q\in\ell}F_{\typ(q),T}(q/n)
                   +2\varepsilon TC_{\rm line}.
\]
For this input the field bound is 7 and there are \(256^2\)
common rectangles. The lattice-to-chord argument without removing
the root has error at most
\(E_{\rm cap}/n\), where \(E_{\rm cap}=4\cdot256^2\cdot7=1835008\).
Take \(\varepsilon=\varepsilon_H\) and
\(n\ge4E_{\rm cap}/\kappa_H\). The lattice error and ledger
error are each at most \(\kappa_H/4\), so \(k_\ell<K\).
Finally \(Z_p(T)\le B\) changes to any point law at cost at most
\(B=n^d\), proving \eqref{eq:clockexplicit}.

To verify the stated threshold, put \(s=\log n\ge10^7\).
The rational bounds \(\varepsilon_H>2^{-34}\), \(T<2^{11}\),
\(\kappa_H>2^{-11}\), and \(E_{\rm cap}<2^{21}\) imply all
three preliminary conditions, since \(n^{1-d}=e^{s/20000}\ge
e^{500}>2^{500}\). The right side of \eqref{eq:clockexplicit}
is at most
\[
 2^{18}\exp\left(2s-\frac{e^{s/40000}}{2^{86}s}\right).
\]
The function \(e^{s/40000}/s^2\) increases for \(s\ge80000\).
At \(s=10^7\), use \(e^{250}>2^{250}\) and
\(10^{14}<2^{47}\) to obtain
\(e^{s/40000}/(2^{86}s)>4s\); this persists for larger \(s\).
The result is bounded by \(2^{18}e^{-2s}<2\cdot10^{-6}\).
This is only a threshold for the clock cause, not a value of the
eventual threshold in the final existence theorem.
\end{proof}

Combining the three guards shows, for every fixed \(A_0'>0\),
\begin{equation}\label{eq:globaltail}
 \Prob(\sigma\le T)\le C_{A_0'}n^{-A_0'}.
\end{equation}
Under any point law this is at most \(B\) times as large, hence is
\(o(1)\) uniformly in the root. This establishes global-stop smallness
before the weight-cap and ordinary-line tag estimates below.

